\section*{Chapter \ref{ch:g7:money} --- Money}

\begin{solution}{exo:g7:money:1}
About $30\,\%$ in 2011 and about $63\,\%$ in 2024 (the table gives $29.7\,\%$ and $63.3\,\%$).
\end{solution}

\begin{solution}{exo:g7:money:2}
(a) Fish--grain, fish--pots, fish--cloth, grain--pots, grain--cloth, pots--cloth: 6 rates.
(b) 4 \omterm{def:g7:money:price}{prices}, one per good.
\end{solution}

\begin{solution}{exo:g7:money:3}
(a) \omterm{def:g7:money:unit-of-account}{Unit of account}. (b) \omterm{def:g7:money:store-of-value}{Store of value}. (c) \omterm{def:g7:money:means-of-payment}{Means of payment}.
\end{solution}

\begin{solution}{exo:g7:money:4}
Bread $2 \div 2 = 1$ loaf; milk $1.5 \div 2 = 0.75$ loaf; cheese $6 \div 2 = 3$ loaves.
\end{solution}

\begin{solution}{exo:g7:money:5}
Each of the 10 goods pairs with the 9 others: $10 \times 9 = 90$; but this counts every pair
twice (fish--grain and grain--fish), so there are $90 \div 2 = 45$ rates.
\end{solution}

\begin{solution}{exo:g7:money:6}
(a) Bread $8 \times 0.25 = 2$ dalens; jam $12 \times 0.25 = 3$ dalens. (b) $12 \div 8 = 1.5$
rations (or $3 \div 2 = 1.5$). (c) No: the unit changes both \omterm{def:g7:money:price}{prices} in the same proportion.
\end{solution}

\begin{solution}{exo:g7:money:7}
(a) $80 \to 40 \to 20 \to 10$ loaves. (b) \omterm{def:g7:money:store-of-value}{Store of value}. (c) Spend it as soon as they get it,
or keep their savings in goods or in another \omterm{def:g7:money:money}{money}.
\end{solution}

\begin{solution}{exo:g7:money:8}
Kenya: $90.1 - 42.3 = 47.8$ percentage points; the world: $78.7 - 50.6 = 28.1$ percentage
points.
\end{solution}

\begin{solution}{exo:g7:money:9}
For example: a large market of strangers trading many goods, where trading is costly, that
keeps bartering for generations without any good coming into use as \omterm{def:g7:money:money}{money}; or a good that
rots within days becoming the \omterm{def:g7:money:money}{money} of a society that had durable goods available.
\end{solution}

\begin{solution}{exo:g7:money:10}
(a) 2 goods: 1 rate; 3: 3; 4: 6; 5: 10; 6: 15. (b) The new good needs a rate with each of
the goods already there: the third good adds 2, the fourth 3, the fifth 4, and so on.
(c) From 10 goods (45 rates), the 11th adds 10 and the 12th adds 11: $45 + 10 + 11 = 66$
rates.
\end{solution}

\begin{solution}{exo:g7:money:11}
(a) It grows: each coin must be checked, a minute per coin, and false coins are refused.
(b) Coins of another country, notes or payments from an account, or even goods (\omterm{def:g7:money:barter}{barter}).
(c) Being hard to forge.
\end{solution}

\begin{solution}{exo:g7:money:12}
For example: how often people pay in cash, by card or by phone; how much \omterm{def:g7:money:money}{money} the accounts
hold; whether the accounts are used at all. A survey in which people write down every
payment they make for some days (a payment diary) would answer the first.
\end{solution}

\begin{solution}{pb:g7:money:1}
\textbf{1.} No: in each pair, at most one trader wants what the other has (for example, Bo
wants rope, not Ana's fish).
\textbf{2.} $5 \times 4 \div 2 = 10$ rates.
\textbf{3.} Ana gives fish to Eli for salt, salt to Di for pots, pots to Cy for rope, and rope
to Bo for bread: 4 swaps.
\textbf{4.} 5 \omterm{def:g7:money:price}{prices}, one per good.
\textbf{5.} She receives 4 shells and pays 2: she has 2 shells left over.
\textbf{6.} A pot is worth $8 \div 2 = 4$ loaves; rope is worth $5 \div 1 = 5$ units of salt.
\textbf{7.} (a) \omterm{def:g7:money:unit-of-account}{Unit of account}. (b) \omterm{def:g7:money:store-of-value}{Store of value}.
\textbf{8.} If anyone could find or copy shells, their number would jump, each shell would
buy less, and nobody could trust what a shell is worth.
\textbf{9.} Every \omterm{def:g7:money:price}{price} doubles: fish 8, bread 4, rope 10, pots 16, salt 2 shells.
\textbf{10.} $16 \div 4 = 4$ loaves: the same as in question 6.
\textbf{11.} Before: $10 \div 2 = 5$ loaves; after: $10 \div 4 = 2.5$ loaves.
\textbf{12.} Those who collected the new shells gain: they can buy goods with them. Those who
had saved shells, like Bo, lose: their shells buy half as much. The goods on the island have
not changed.
\end{solution}
